
\resumenCastellano{
Este trabajo aborda la heterogeneidad existente en los programas de Ingeniería Informática en Chile mediante un enfoque de análisis computacional de mallas curriculares y descripciones de asignaturas. 
Se propone una metodología basada en procesamiento de lenguaje natural para la representación textual de los contenidos formativos, junto con técnicas de aprendizaje no supervisado, específicamente clustering, 
para identificar patrones estructurales entre programas. Adicionalmente, se incorporan métricas de similitud y análisis comparativo de perfiles de egreso y distribución de áreas disciplinares, con el fin de complementar el agrupamiento y proporcionar una caracterización más integral. 
El estudio considera datos recopilados desde fuentes institucionales públicas, permitiendo construir una visión sistemática y actualizada de la oferta académica. 
Como resultado, se espera evidenciar diferencias significativas en enfoques formativos, niveles de especialización y coherencia curricular entre instituciones, contribuyendo a una mejor comprensión del sistema y proporcionando una base objetiva para la comparación de programas.

\vspace*{0.5cm}
\KeywordsES{Educación en Ingeniería; Algoritmos de agrupamiento; Procesamiento de lenguaje natural; Programas educativos; Minería de textos}
}